\documentclass[instructions.tex]{subfiles}
\begin{document}
 
\chapter{Faux raisonnement sur le salut}
 

\primo Personne ici-bas ne sait s'il est prédestiné, s'il est un vase de miséricorde qui sera élu, ou un vase de colère qui sera brisé\footnote{Nescit homo utrum amore an odio dignus sit. Eccl.9.}. C'est un secret réservé à Dieu seul. Dans cette cruelle incertitude, l'homme sage travaille à son salut avec crainte, dans la confiance que Dieu lui accordera sa miséricorde. 

Mais voici le ridicule raisonnement dont l'esprit de ténèbres amuse les mondains: \og Ou je suis prédestiné ou je ne le suis pas. Si je le suis, quelques crimes que je commette, je serai sauvé. Si je ne suis pas prédestiné, je serai donc, quoi que je fasse, infailliblement damné\fg; ainsi, concluent-ils, \og il est indifférent de faire le bien ou le mal pour être sauvé\fg. Ce raisonnement est affreux; il faut avoir perdu le bon sens pour fonder sa conduite sur un raisonnement si frivole. 

Si Dieu prédestine les uns, s'il réprouve les autres, c'est parce qu'ils auront bien ou mal vécu. Ainsi, si vous êtes un jour réprouvé, c'est parce que vous le mériterez par votre vie et par votre mort dans le péché mortel. Vivez donc de manière que vous puissiez mourir en grâce, et vous ne serez point réprouvé. De même si vous êtes prédestiné, c'est parce que Dieu vous accordera de mourir dans sa grâce après avoir mené une vie sainte. Vivez donc dans la sainteté, pour mourir saintement, et vous serez prédestiné. 

Quoique Dieu veuille sauver tous les hommes, sa volonté seule ne rend pas les hommes saints. Il ne veut donner de gloire qu'à ceux qui auront vécu dans l'innocence ou la pénitence, de même qu'il ne veut condamner à l'enfer que ceux qui auront vécu ou qui seront morts dans le péché. S'il vous y condamne, c'est parce que vous le mériterez, et non pas précisément parce qu'il le veut. 

Lorsqu'un juge condamne un criminel selon la loi, ce n'est ni le juge ni la loi précisément qui sont cause de la condamnation, mais le crime qui a été commis; de même ce n'est pas précisément parce qu'il lui plait, que Dieu condamne à l'enfer, c'est parce qu'on le mérite. Un juge qui prononce la sentence ne dit pas: \og Je te condamne parce que je le veux\fg; mais: \og Je te condamne parce que tu le mérites\fg. Raisonnons de même des jugements de Dieu. 
Quoiqu'il sache ce que nous ferons, cette prévision de Dieu n'ôte rien à notre liberté, parce que la connaissance de Dieu n'est point la cause des événements. Les choses arrivent comme Dieu les a prévues, \og mais elles n'arrivent pas parce qu'il les a prévues: au contraire, il les a prévues et il les connait, parce qu'elles arriveront\fg, dit saint Augustin. Quand je vois un voyageur qui s'égare, ce n'est point ma connaissance qui est cause de son égarement, mais l'ignorance de ce voyageur, qui ne sait pas la route. De même, quand Dieu prévoit notre égarement et notre malheur éternel, ce n'est point à lui, mais à nous, qu'il faut l'attribuer. 

\secundo Le salut dépend de Dieu, qui nous appelle tous et qui nous aide tous par sa grâce; mais il dépend aussi de notre coopération. Ainsi, quand vous sauriez par révélation que vous serez sauvé, vous ne devriez pas pour cela cesser de bien vivre. Un laboureur cesse-t-il de cultiver la terre parce qu'il connait que la récolte sera abondante ? N'est-ce pas parce qu'il espère la récolte, qu'il se détermine à semer ? 

Oh! que l'esprit de l'homme est bizarre quand il se livre à ses égaremens ! On dirait que certaines gens n'ont de l'esprit que pour s'aveugler sur le salut, tandis qu'ils raisonnent si prudemment sur les affaires du monde. Que répondriez-vous à celui qui dirait: Ou Dieu voit que je moissonnerai, ou il voit que je ne moissonnerai rien; s'il voit que je moissonnerai, quoi qu'il arrive, je ferai la moisson; ainsi je ne dois me mettre en peine de semer ? Raisonnement extravagant, parce qu'il voit que vous aurez semé un grain qui poussera un germe, et qui croîtra. Il voit que vous ne moissonnerez pas, si vous ne jetez pas de semence en terre. Raisonnons de même sur le salut. Semez dans le temps et vous moissonnerez dans l'éternité\dots

\og Celui qui vous a fait sans vous\fg, dit saint Augustin, \og ne vous sauvera pas sans vous\fg. C'est pour cela que le démon tâche par ses tentations de nous perdre, en nous empêchant de coopérer à la grâce et aux desseins de Dieu, parce qu'il sait que, si notre prédestination dépend principalement de la grâce, il faut aussi notre coopération. Si Dieu ne nous sauve pas sans nous, le démon ne peut aussi nous perdre sans nous. S'il tâche de nous perdre, il ne tient qu'à nous de rendre ses efforts inutiles. C'est un chien furieux, qui peut aboyer, dis saint Augustin, mais qui ne peut mordre que ceux qui le veulent\footnote{Latrare potest, mordere non potest, nisi volentem. S.Aug.}. Ce n'est donc pas à la volonté de Dieu ni à la malice du démon, que nous devons attribuer notre perte, mais à notre négligence et à notre malice. 

Voulez-vous être sauvé, faites ce que vous feriez si vous étiez assuré de l'être. Tâchez, dit saint Pierre, d'assurer votre prédestination par les bonnes oeuvres\footnote{Satagite ut per bona opera certam, vestram vocationem et electionem faciatis ? Petr.1.10}. C'est en vain qu'on raisonne, qu'on pointille sur le mystère de la prédestination, il faudra toujours en venir à ce point:\textit{On ne moissonera que ce que l'on aura semé}. 

\section{Résolutions}

\primo J'espèrerai toujours de pouvoir réussir dans la grande affaire de mon salut, aidé par la grâce divine. 

\secundo Mais aussi je ferai tout ce qui dépendra de moi, pour correspondre aux grâces que le Seigneur daignera m'accorder dans sa miséricorde.

O mon Dieu! Je sais que la persévérance dans votre amour est un don spécial de votre grâce; accordez-moi ce grand don, que je vous demanderai tous les jours de ma vie avec confiance.

\end{document}
